Read in Gauge Theory Gravity: the metric $g$ below is the derived metric of the gravitational field $(\hb, \Omega)$ on flat spacetime, the manifold is the invertible region, and the tractor connection is the flat-space $\Cl(2,4)$ gauge field of Theorem \claimref{synthesis.gtg-tractor-gauge-field-reproduces-tractor-connection}; the index computation below is its reading in the splitting of $g$ (Fact \claimref{synthesis.imported-gtg-dictionary}).

Constraints: $\delta Q(I) = 2\ip{\delta I}{I} = 0$, $\delta\ip{V}{I} = \ip{\delta V}{I} + \ip{V}{\delta I} = 0$, $\delta Q(V) = 2\ip{\delta V}{V} = 0$: three independent linear conditions on the twelve components of $(\delta I, \delta V)$ (independent because $I$ and $V$ are non-null and orthogonal), leaving nine. Stabiliser: a bivector annihilating both $I$ and $V$ is a bivector of $\{I,V\}^\perp$, which has signature $(1,3)$ (Theorem \claimref{dark-energy.frozen-point-makes-stelle-west-vector}(ii)), so the stabiliser algebra is $\mathfrak{so}(1,3)$, of dimension six, and the broken complement has dimension nine. The map $\xi \mapsto (\xi\cdot I, \xi\cdot V)$ maps into the constrained space, since $\ip{\xi\cdot I}{I} = 0$, $\ip{\xi\cdot V}{V} = 0$ and $\ip{\xi\cdot V}{I} + \ip{V}{\xi\cdot I} = 0$ by antisymmetry of $\xi$; restricted to the broken complement it is injective, because a $\xi$ in the complement with $\xi\cdot I = \xi\cdot V = 0$ lies in the stabiliser; by equal dimensions it is bijective. Any constrained $(\delta I, \delta V)$ is thus $(\xi\cdot I, \xi\cdot V)$ for a unique broken $\xi$, and the gauge transformation by $-\xi$, acting simultaneously on $\nabla$, $I$ and $V$, removes it while shifting $\delta\nabla$: the Goldstone eating. In the resulting gauge the second variation is a quadratic form in $\delta\nabla$ alone. Invariance of the functional: a gauge transformation is a section $g$ of the ambient group bundle, acting on the fibre by $g(x) \in SO(2,4)$ at each point; it acts on the connection by $\nabla \mapsto g\nabla g^{-1}$, so that the curvature, $F = \nabla^2$, transforms by $F \mapsto gFg^{-1}$, that is, its bivector value at each point is carried by $g(x)$ (a bivector $a\wedge b$ goes to $ga\wedge gb$, since $g$ is an isometry), and on $I$ and $V$ by $I \mapsto gI$, $V \mapsto gV$. The volume form of a six-dimensional fibre with a nondegenerate form is invariant under its special orthogonal group, so at each point $\epsilon(gI, gV, gF, gF) = \det(g)\,\epsilon(I, V, F, F) = \epsilon(I, V, F, F)$; the prefactor $1/Q(I)$ and the three constraints $Q(I) = \Lambda/3$, $\ip{V}{I} = 0$, $Q(V) = -3/\Lambda$ are preserved because $g(x)$ is an isometry of the fibre; and a pointwise identity between integrands is an identity between the integrals. The functional is therefore invariant, and the constraint surface is carried to itself, under the ambient gauge group acting on all three arguments. The script confirms the fibrewise identity on rational elements of $SO(2,4)$ (a boost, a rotation, and their product, each confirmed an isometry). That the one-argument version fails is worth recording, since it is the error to avoid: with $\xi = I\wedge V$ one has $\xi\cdot I = -(\Lambda/3)V$, a tangential fluctuation of $I$, but $\xi\cdot X = -(3/\Lambda)I - V$, so that $\epsilon(I, \xi\cdot X, F, F) = -\epsilon(I, V, F, F)$: acting on $\nabla$ and $I$ while holding $X$ fixed changes the integrand by minus itself, not by zero. The script performs the counting exactly in the flat model over $\mathbb{Q}$ at $\Lambda = 3$: the joint action matrix of the fifteen basis bivectors on $(I, V)$ has rank exactly nine, so the stabiliser has dimension six; the image of the bivector algebra on $I$ alone has rank five and equals $I^\perp$; and the one-argument failure is confirmed on the displayed $\xi$.
